\subsection{Two-mark confluence and angular-to-lattice scales}\label{c4:pair-confluence}

The conjugation gives two angular expressions for the same trace. We
first identify its confluent value with the finite color numerator, then
compare its self factors with those of the radial transform. The general
oscillator transformation is already available from
Section~\ref{c4:radial}; the next subsection will derive the slope and
phase rules for the present charges. Keeping the self factors here makes
their cancellation against the Pfaffian normalization explicit.

\paragraph{Fusing the site cluster.} Use \(\mathcal L^*\) for all momenta, cocycle as above
\[
 \exp(\pi i[b^t J P/2-3 k_b(p_C,P)]),\quad p_C=\Pi_C h_1,\quad k_b\text{ the external }v\text{-label coefficient}.
\]
Thus \(J\) has first row \((0,1,1,0,0,0,-1,-1)\) in the first sector, opposite in second, zero across. Color fields are the combinations with coefficients \(1,q^{-3},q^{-6}\) on the angular side, of charges \(h_i,h_i-\alpha_{i0},h_i-\alpha_{i0}-\alpha_{i1}\), denoted \(b_{ic}\) with \(c=-1,0,1\). The two mark insertions have charge \(v\); add the compensator of charge \(-2v\) first, at \(s_*=u_*\) real positive. Work with \(0<u_*<L\). Use unshifted and shifted cluster positions \(s_j=2\pi i e_j+\epsilon\lambda_j\) (\(e_j=0,1\)), taking unshifted angles slightly positive and shifted slightly less than \(2\pi\), all distinct. The corresponding lattice site parameters are proportional to \(\exp(-s_j/8)\); in the limit their ratio between classes is \(q^{\pm1}\).

The two marked sites have a common shift \(e\in\{0,1\}\), as required
in the two-mark bound. Write \(\mathcal V_j(s_j)\) for the color combination or mark field just
specified, evaluated at \(z(s_j)=e^{-2\pi i s_j/L}\) and scaled field by
field by \((2\pi/L)^{b^2/2}\). Let \(\mathcal V_*(u_*)\) be the similarly
scaled compensator. The normalized angular trace is
\[
 \mathcal T_\epsilon(L,u_*)=
 \frac{\Tr\!\left[\widehat C Q_L^H
       \left(\prod_j^{\leftarrow\mathrm{ang}}\mathcal V_j(s_j)\right)
       \mathcal V_*(u_*)\right]}
      {Z_C\Theta_C^{\mathrm{ang}}(L)},\qquad
 \Theta_C^{\mathrm{ang}}(L)=
       \sum_{t\in\mathbb Z}Q_L^{(4tp_C)^2/2}.
\]
The arrow orders the fields latest first in angle. We use the
coefficientwise trace convention of Section~\ref{c4:angular}, and then
the continuation justified below.

Moving the shifted (latest) group cyclically over the twist brings them to positions \(s'_j=s_j-2\pi i e_j\), earliest, by \(Q_L^H W(Q_L^{-1}z)=W(z)Q_L^H\), with charges and fields now acted on by \(\widehat C\). There is no extra cocycle because \(J,p_C\) are preserved. Conjugate then as above. Thus the pure cluster charges are \(u_j=A^{e_j}h_i\) (or \(A^{e_j}v\) at the marks). With \(j\prec k\) denoting their fixed post-move chronological order, define the divided trace
\begin{equation}\label{c4:confluent-trace}
 \mathcal T(L,u_*)=
 \lim_{\epsilon\to0^+}
 \frac{\mathcal T_\epsilon(L,u_*)}
      {\prod_{j\prec k}(s'_j-s'_k)^{(u_j,u_k)}}.
\end{equation}
Fix the branches throughout. On the conjugated side the cluster fuses
to its total charge. The division supplies exactly the vertex scaling
as for \eqref{c4:Ang-2}; cluster cocycles contribute only a fixed unit
factor.

As before all formal equalities are first read at fixed \(H\) and field powers (with the spread estimates for finiteness during cut changes), then evaluated near small \(Q_L\) by angular Gaussian and absolute oscillator bounds. The ordering for the post-conjugation root contours can put the later group on the far side of the external positions by commuting, as described in the ray estimate. Compensator-to-shifted angular separation in the original \(C\)-side tends to an entire turn, but \(0<u_*<L\) separates the points for continuation near confluence. Thus the \(C\)-side has a divided Laurent expansion with at worst finite pole, real-analytic coefficients at all positive \(L\) (keeping \(u_*/L\) fixed); negative coefficients vanish by the small-\(Q_L\) equality. The fused-vertex continuation follows by the analytic bound \eqref{c4:Spread-3} on the radial series.

\paragraph{The angular limit.} For each fixed finite site array, we evaluate \(|\mathcal T(L,u_*)|\) by first letting \(L\to\infty\) at fixed \(u_*\), then \(u_*\to\infty\). The array size will grow only after these limits. Terms require \(\Pi_C d=0,\ d=\sum b_j\) now including compensator, i.e. equal color sums in both sectors since \((p_C,b_{ic})=(-1)^i c/4,\ (p_C,v)=0,\ p_C^2=1/4\). Conversely every such term has a momentum solution \(m=n+d,\ n=Cm,\ m=-K_C d+t p_C,\ K_C=(C-I)^{-1}\) off the invariant line, with a full arithmetic progression of \(t\) of spacing 4. Indeed the within-orbit differences of pairings \((\alpha_{is},K_Cd)\) are integers by balance and \(d\in\mathcal L^*\); also \((\alpha_{10}+\alpha_{20},K_C b_{ic})\) is integral (multiply the eight shifts). The two marks cancel the compensator in \(d\). All terms using \(J,k_b\) are independent of \(t\); the heat and angular-power sum divided by the invariant partition tends to 1 at bounded \(u\), uniformly locally. The products converge as in \eqref{c4:Ang-3}, now to pair products of \(P_l(\tau)=2\sin(\pi(l+\tau)/8)\) with \(\tau=i(s_j-s_k)/(2\pi)\) for chronological \(j<k\), raised to powers \((C^l b_j,b_k)\), and scaled self factors \(8^{-b_j^2/2}\prod_{l=1}^7 P_l(0)^{(C^l b_j,b_j)/2}\). Although the exponents can have nonzero mean, constants cancel globally by \((\Pi_C d)^2=0\) (subtract the same reference log-product for offsets \(8m+4,8m+4\) for each \(l\)). Branches are positive real for real \(0<\tau<1\). The limits commute with Laurent coefficient extraction by the same groupwise uniform bounds after factoring finitely many singular powers. The compensator pairs then tend to unit phases \(\exp(2\pi i(K_Cv,b))\) since its invariant projection is zero and \(\Im\tau\to+\infty\).

\paragraph{Pair powers and phases.} For explicit comparison with the colored sine products write
\[
 H=(4,2,-2,-4,-2,-4,-2,2)/3,\qquad H'=(4,2,-2,-4,-2,2)/3 .
\]
Subtracting \(H_l\) from the \(C\)-pair exponents gives exactly the integral color exponents \(\ell_l\) above (multiply the circulant shifts); for two marks instead the difference is \(1/4\) constant. The signs for these integral factors in chronological angular order are \((-1)^{\sigma_{ij}}\) with \(\sigma=1\) on the directed within-sector cycle \(-\to0\to+\to-\), zero otherwise (for an external mark paired with an occupied color also a minus, canceling since there are two marks). This description works even if the color product was first given in another order: reversal negates \(P_0\) with parity exactly swapping the signs on the cycle and reflects the other indices without sign.

In units \(\pi i\), the \(J\)-phase is minus the ordered-pair sum of \(b_j^t J K_C C b_k\), including half of the self terms, by \(n=-K_C C d\) off the invariant line and skew symmetry. Add the correction \(-3\sum_{j<k} k_{b_k}(p_C,b_j)\). The net from these is precisely the required sign up to global phases. Here are checks: putting \(M=JK_CC\) the cross-sector color table is zero, and subtracting the central entry and \(D_i c c'\) in \(b_{ic}^t M b_{ic'}+\sigma_{cc'}\) leaves the tables
\[
 \begin{pmatrix}0&0&-4\\0&0&0\\0&0&0\end{pmatrix},\qquad
 \begin{pmatrix}2&2&2\\0&0&2\\0&0&2\end{pmatrix},\qquad D_1=-9/4,\quad D_2=1/4.
\]
Thus the common-sum squares contribute \(-(\sum_{\text{sector }1} c)^2\) modulo 2 (sum here only over first-sector sites), which cancels the half-self parity in sector 2. For the external terms one may move \(v,v,-2v\) formally to the front of this phase sum since \(b_{ic}^t J v+3(p_C,b_{ic})=0\); their linear contributions cancel. Finally \((\alpha_{i0},K_Cv)=(\alpha_{i1},K_Cv)=-3/8\); the separated-compensator sine phases thus cancel color dependence of the local angular coefficients above. Relative to the scaled baseline self factor using \(H\), color self factors have ratio \(P_2(0)P_3(0)\) at \(\pm\), \(P_2(0)P_3(0)/P_1(0)\) at zero. This is \(2+\sqrt2\) times the local weights \((\sqrt{2-\sqrt2},1,\sqrt{2-\sqrt2})\).

\paragraph{The normalized comparison.} For a period-\(d'\) array \(D\), set
\[
 M(D,r)= d'^{-D_{-r}}
    \prod_{\substack{0\le l<d'\\l\ne-r\bmod d'}}|1-e^{2\pi i(l+r)/d'}|^{D_l}.
\]
The baseline \(C\)-sine powers after cluster division give \(M(H,e_j-e_k)\) per pair and \(M(H,0)^{1/2}\) per site; indeed the zero sine divided by the merging coordinate difference tends in modulus to \(1/8\). Two marks with their extra \(1/4\) powers only change by a bounded positive factor (divide there with the appropriate pure exponents \((u_j,u_k)\)); similarly for mark self factors. Consequently
\begin{equation}\label{c4:confluent-color-scale}
 \lim_{u_*\to\infty}\lim_{L\to\infty}|\mathcal T(L,u_*)|
 \asymp |\mathscr N|(2+\sqrt2)^N M(H,0)^{N/2}
                  \prod_{j<k} M(H,e_j-e_k).
\end{equation}
The comparison factors are bounded above and below independently of
the site counts and shift proportions, including both common shifts of
the two marks.
Here color normalization uses the displayed local weights and the integral pair factors. Possible singular individual color terms are interpreted by confluence: the full rational lattice sum has a regular limit and common baseline factors are extracted before specialization.

The radial transform will use the heavy self scale
\begin{equation}\label{c4:heavy-self-scale}
 \mathcal S=M(H',0)^{N/2}\prod_{j<k}M(H',e_j-e_k)>0.
\end{equation}
Indeed \((A^k h_i,h_j)=(A^k h_i,v)=H'_k\) and
\(A^2 h_i-Ah_i+h_i=0\), so the six-offset self calculation of
\eqref{c4:Rad-1} applies to every term involving \(h_i\); the remaining
two-mark-only factors are bounded independently of the counts.
Explicitly
\[\begin{gathered}
 M(H',0)/M(H,0)=[(2+\sqrt2)/(2b')]^2,\\
 M(H',1)/M(H,1)=[2\sqrt2/(\sqrt3\sqrt{2-\sqrt2})]^2,\qquad b'=3/4.
\end{gathered}\]
These are \(g(0)^2,g(\pi/4)^2\) of \eqref{c4:Pair-Pf}.
Dividing \eqref{c4:confluent-color-scale} by \(\mathcal S\) therefore
leaves, within bounded factors,
\[
 |\mathscr N|(2b')^N\prod_{j<k}g(\theta_j-\theta_k)^{-2},
 \qquad \theta_j-\theta_k\in\{0,\pm\pi/4\}.
\]
Since \(\mathscr H=\mathscr N/\widetilde D^2\), the Pfaffian lower
bound gives the precise comparison used below:
\begin{equation}\label{c4:normalized-two-mark-comparison}
 |\mathscr H|\le C N^{-5/24+o(1)}
       \lim_{u_*\to\infty}\lim_{L\to\infty}
                       \frac{|\mathcal T(L,u_*)|}{\mathcal S}.
\end{equation}
The constant and exponent error are uniform in the shift proportions.
The next subsection supplies the radial expression for this same
\(\mathcal T\). Section~\ref{c4:hierarchical} bounds its absolute value
uniformly through the two successive limits before letting the array
sizes grow.
